% ---------------------------------------------------------
\section{Approximating Limits via Difference Quotients}

\begin{propositionbox}{Definition (Difference Quotient)}
The \emph{difference quotient} of a function $f$ at $x = a$ with
increment $h \neq 0$ is
\[
  \frac{f(a + h) - f(a)}{h}.
\]
\end{propositionbox}

\begin{remark*}[Geometric meaning]
The difference quotient is the slope of the secant line through
the points $(a,\, f(a))$ and $(a+h,\, f(a+h))$ on the graph of $f$.
As $h \to 0$, the secant line rotates toward the tangent line at
$(a, f(a))$, and the difference quotient approaches the slope of
that tangent --- provided the limit exists.
\end{remark*}

\begin{remark*}[Approximating a limit numerically]
For a function whose limit at $a$ is difficult to evaluate
algebraically, the difference quotient with small values of $h$
provides a numerical estimate.
Evaluating at $h = 0.1,\, 0.01,\, 0.001$ and observing the
output stabilise gives evidence for the limiting value.
For example, to estimate $\lim_{x \to 0} \dfrac{\sin x}{x}$:

\medskip
\begin{center}
\begin{tabular}{@{}cc@{}}
\textbf{$h$} & \textbf{$\sin(h)/h$} \\
\hline
$0.1$    & $0.99833$ \\
$0.01$   & $0.99998$ \\
$0.001$  & $0.99999\overline{9}$ \\
\end{tabular}
\end{center}
\medskip

The values converge to $1$, consistent with the exact result
$\lim_{x \to 0} \dfrac{\sin x}{x} = 1$.
\end{remark*}

\begin{remark*}[Limitation of numerical approximation]
Numerical evidence is not a proof.
Taking $h$ small can introduce floating-point cancellation error,
and a sequence of outputs that appears to stabilise may still
diverge or converge to the wrong value for poorly behaved
functions.
The difference quotient is a computational tool for building
intuition and checking algebra; the limit itself requires the
$\varepsilon$-$\delta$ definition for a rigorous conclusion.
\end{remark*}